\chapter{Avaliação e Discussão}
\label{cap:avaliacao}

\comment[inline]{Capítulo reflexivo: avaliar o que foi construído, confrontar com os objetivos declarados, discutir limitações e implicações. Aproximadamente 6--8 páginas.}

\section{Avaliação da Usabilidade}

\comment[inline]{Quinto objetivo específico: comparar as visualizações produzidas com o formato original do Portal da Transparência. Definir critério de comparação: o que é possível responder com o DW que não era possível antes? Usar exemplos concretos de perguntas e mostrar como cada sistema as responde (ou não). Pode ser uma tabela comparativa de tarefas.}

\comment[inline]{Se houver avaliação com usuários (questionário, entrevista, teste de usabilidade): descrever a metodologia, o perfil dos participantes e os resultados. Mesmo um piloto pequeno (5--8 pessoas) adiciona valor empírico ao trabalho.}

\section{Limitações do Trabalho}

\comment[inline]{Ser honesto sobre as limitações: cobertura temporal dos dados, qualidade das fontes, funcionalidades de BI não implementadas, ausência de atualização automática, restrições da ferramenta de visualização escolhida. Limitações bem documentadas demonstram maturidade científica.}

\section{Implicações para Transparência e Controle Social}

\comment[inline]{Discussão teórica: em que medida o DW construído endereça as três lacunas identificadas na introdução (social, técnica, pedagógica)? Retomar Freire e Kimball como fios condutores. O que este trabalho diz sobre o papel da arquitetura de dados na democracia? Esta seção fecha o argumento iniciado na introdução.}

\section{Replicabilidade e Escalabilidade}

\comment[inline]{O modelo é replicável para outros municípios? Quais adaptações seriam necessárias? Discutir o potencial de escala da abordagem: o que seria necessário para aplicar este modelo a outros 20 municípios mineiros, por exemplo? Isso valoriza a contribuição metodológica do trabalho.}
